\section{Discussion}
\label{sec:discussion}

\subsection{Relationship to the broader space-MBSE review}
\label{sec:discussion-buruso-relationship}

Buruso et al. provide the appropriate background for this review because their comprehensive space-MBSE synthesis already covers the broad distribution of MBSE methods, modeling languages, tools, lifecycle phases, and enabler classes across space applications \cite{Buruso2026MBSESpace}.  The present review therefore does not treat the recurrence of SysML, Cameo/MagicDraw, Capella/Arcadia, PLM tools, simulation environments, or integrated digital-engineering toolchains as a novel finding by itself.  Instead, those enablers are interpreted only where they affect nanosatellite and EPS-relevant value evidence, implementation burden, or model-linked V\&V.

The strongest connection to Buruso et al. is their discussion of value-driven implementation and their Table~5 research gaps.  They identify persistent gaps in practical implementation pathways, fragmented toolchains and lifecycle interoperability, resource-mismatched frameworks for small missions, weak linkage between MBSE artifacts, governance and ROI, and misalignment with NewSpace-style development needs \cite{Buruso2026MBSESpace}.  Our results support those gaps in a narrower population and make their nanosatellite consequences more explicit.  In the included studies, practical implementation pathways and toolchain interoperability are frequently addressed, with evidence-backed classifications in 29 and 19 studies respectively, but governance/ROI/value and NewSpace/agile/distributed development are less frequently addressed with evidence (7 and 6 studies respectively; Table~\ref{tab:results-gap-mapping-counts}).  Thus, the nanosatellite literature is not simply missing MBSE awareness; it is missing durable evidence that the models change project outcomes, review decisions, or downstream verification practice.

This narrower focus also exposes an EPS-specific asymmetry that is hidden in broad space-MBSE taxonomies.  Many studies include EPS or power concerns inside whole-CubeSat models, but only two included studies are EPS-specific or EPS-first \cite{paper-0026,paper-0059}.  Most examples of power modeling are embedded in system-level trade studies, reference-model work, or preliminary design workflows such as RAX energy/data simulation \cite{paper-0001,paper-0003}, SysML power-budget margins \cite{paper-0080}, Capella/Arcadia preliminary EPS allocation \cite{paper-0078}, PLM/Dymola battery simulation \cite{paper-0059}, and electrical-architecture concurrent engineering \cite{paper-0053}.  This extends Buruso et al.'s Table~5 concern about implementation pathways: for nanosatellites, the missing pathway is not only ``how to adopt MBSE,'' but how to maintain an EPS/system thread from mission objectives to power requirements, architectural allocation, analysis, verification procedure, and test evidence without imposing enterprise-scale overhead.

\subsection{Value-driven MBSE remains under-evidenced for nanosatellites}
\label{sec:discussion-value-evidence}

The review supports Buruso et al.'s claim that MBSE value must be tied to explicit objectives, model usage, governance, and evaluation criteria \cite{Buruso2026MBSESpace}.  The nanosatellite literature reviewed here contains many value claims, but the evidence is generally modest.  RQ2 was coded as high for only 3 studies and partial for 37, and among 270 extracted value-claim records only 2 were classified as strong evidence; 77 were moderate, 139 weak, and 52 unsupported (Table~\ref{tab:results-study-quality-summary}).  The most common domains were traceability, interoperability/lifecycle continuity, reproducibility/reuse, reliability/mission assurance, and cost/schedule/automation or decision-support mechanisms, but the strong evidence was confined to cost/schedule claims in one concurrent-design comparison rather than EPS reliability or verification closure (Table~\ref{tab:results-claim-evidence-by-domain}).

Moderate evidence usually came from executable demonstrations or traceable artifacts rather than from measured project outcomes.  For example, the RAX studies show how an integrated SysML--ModelCenter--STK--MATLAB workflow can simulate energy and data interactions, support solar-panel and battery-capacity trades, and check mission requirements \cite{paper-0001,paper-0003}.  The open-source Papyrus power-subsystem study demonstrates parameter-impact traceability and requirement-constraint flags from model data \cite{paper-0026}.  PLM/Dymola requirements-in-the-loop simulation reports battery state-of-charge behavior against an SoC requirement \cite{paper-0059}, and the electrical-architecture workflow reports a worst-case battery state-of-charge margin for survival/no-illumination conditions \cite{paper-0053}.  These are useful forms of engineering evidence, but they are mostly prototype, single-case, or pre-physical demonstrations.

Evidence becomes thinner when claims move from technical feasibility to project value.  Reported schedule or cost benefits exist, such as a one-month versus three-month feasibility-study comparison and lower estimated mission cost in a concurrent-design case \cite{paper-0061}, or a four-month preliminary-design cycle with Arcadia and weekly sprints for 3ColStar \cite{paper-0078}.  However, such claims are case-specific and rarely supported by controlled baselines, repeated use, independent stakeholder assessment, or lifecycle follow-through.  Similarly, reference models and templates promise reuse and consistency \cite{paper-0011,paper-0015,paper-0077}, while docs-as-code and open-source workflows improve accessibility and reproducibility \cite{paper-0037,paper-0029}, but the included studies seldom demonstrate reuse across independent missions or quantify the effort required to maintain the model.

The implication is that nanosatellite MBSE remains more artifact-driven than value-driven.  Studies frequently show that a model can contain requirements, architectures, budgets, interfaces, simulations, or generated documents.  Fewer studies define the mission/project objective that justifies the modeling effort, the metric used to evaluate success, the baseline practice being improved, or the review gate at which the model changes a decision.  This is especially important for EPS: power budgets, SoC traces, and margin checks are often convincing local analyses, but they are not always connected to explicit mission assurance objectives such as survival probability, operational availability, fault tolerance, or verification closure.  In this respect, our synthesis directly supports Buruso et al.'s Table~5 gap on weak integration between MBSE artifacts, governance, ROI, and value realization, and shows that the gap persists even in a domain where resource constraints should make value justification central \cite{Buruso2026MBSESpace}.

\subsection{V\&V linkage is present but rarely continuous}
\label{sec:discussion-vv-continuity}

The included studies do not show an absence of V\&V linkage; rather, they show many partial linkages that rarely form a continuous model-to-test thread.  EPS/power-specific pre-physical verification appeared in 25 studies, parametric/budget/constraint checks in 23, executable simulation or scenario analysis in 22, test cases/procedures/V\&V activity packages in 21, requirements-to-model traceability or verification matrices in 19, and external tool or data-flow coupling in 11 (Table~\ref{tab:results-vv-linkage-counts}).  These counts indicate that nanosatellite MBSE studies often recognize verification as a modeling concern.  They also align with Buruso et al.'s observation that later lifecycle testing, integration, and verification have poorer visibility in the wider space-MBSE literature \cite{Buruso2026MBSESpace}.

The most common V\&V pattern is requirements traceability plus a local analytical check.  The EPS-specific Papyrus study extracts requirement constraints and power-budget parameters to generate respected/not-respected flags \cite{paper-0026}.  The academic nanosatellite framework links requirements to ``Verified By'' test cases and behavior diagrams such as bus activation and LEOP flow \cite{paper-0023}.  The 3ColStar study links continuous-power requirements to EPS components and mode-level power analyses \cite{paper-0078}.  The NSS-1 SysML case uses verify relationships and parametric constraints to compute positive mass and power margins \cite{paper-0080}.  These examples show useful pre-physical verification, but the verification often stops at traceability, calculation, or planned test artifacts.

Executable simulation provides a second important but still partial pattern.  RAX-style workflows connect SysML behavior and parametrics to STK, MATLAB, and ModelCenter to check energy, data, and mission-scenario behavior \cite{paper-0001,paper-0003}.  Other studies link model data to power/data-budget simulation or mission-mode simulation \cite{paper-0027,paper-0029,paper-0033,paper-0043}, PLM/Dymola battery behavior \cite{paper-0059}, and electrical-architecture energy-balance simulation \cite{paper-0053}.  These studies strengthen pre-build confidence and support design trades, but they do not generally show continuity from the same model elements to formal test procedures, hardware-in-the-loop execution, integration-test data, anomaly handling, or acceptance evidence.

The few test-oriented examples reinforce this distinction.  Requirements Modeling Language and RMM work links requirements, code, tests, and defects for CubeSat software-oriented examples \cite{paper-0034}.  EIRSAT-1 demonstrates a communications-subsystem verification model coupled through Python to ground-station hardware and requirements-compliance updates \cite{paper-0063}.  These are among the clearest model-to-test patterns in the included studies, but they are not EPS demonstrations.  For EPS, the literature more often contains power budgets, margin calculations, and simulations than a continuous thread from EPS requirement to model allocation, analysis case, test case, test procedure, measured test data, and verification closure.  This is the central V\&V gap for nanosatellite MBSE: linkage exists, but it is fragmented across requirements traceability, parametric checks, simulation, and planned V\&V packages.

\subsection{Implications for lightweight nanosatellite MBSE practice}
\label{sec:discussion-practice-implications}

For resource-constrained CubeSat and nanosatellite teams, the practical implication is not to adopt more MBSE enablers, but to adopt a thinner value-and-verification thread.  Buruso et al.'s Table~5 calls for modular, scalable MBSE frameworks and adoption playbooks for lightweight missions \cite{Buruso2026MBSESpace}.  Our synthesis suggests that such playbooks should start by defining the value hypothesis for the model: which decision, review, risk, or verification activity the model is expected to improve; what artifact will be produced; what metric or comparator will indicate success; and when the modeling effort should stop.  Without these criteria, teams can accumulate model artifacts while still lacking evidence that MBSE improved reliability, reproducibility, integrability, or schedule performance.

A lightweight nanosatellite playbook should make EPS a first-class model thread rather than a subsystem table appended to the architecture.  At minimum, the thread should connect mission modes and operational scenarios to EPS requirements, power-generation/storage/distribution architecture, assumptions, margins, budget equations, executable checks, verification methods, and test cases.  Studies already provide pieces of this thread: open-source Papyrus/SysML traceability for a power subsystem \cite{paper-0026}, model-derived power and data budgets \cite{paper-0029}, Capella/PVMT simulation coupling \cite{paper-0043}, PLM/Dymola requirements-in-the-loop behavior \cite{paper-0059}, and concurrent electrical-architecture simulation \cite{paper-0053}.  The missing practice guidance is how a small team should select a minimal subset of these mechanisms and maintain it through review and test readiness.

Tool choices should be evaluated against adoption burden as well as analytical capability.  Enterprise or integrated environments can provide powerful lifecycle continuity, as shown by PLM/Dymola and concurrent-engineering workflows \cite{paper-0059,paper-0053}, but they may reproduce the resource-mismatch problem noted by Buruso et al. \cite{Buruso2026MBSESpace}.  Conversely, open-source, plain-text, or docs-as-code approaches can reduce licensing and training barriers while supporting traceability and generated deliverables \cite{paper-0026,paper-0029,paper-0037}.  The OPAL user study suggests that low-training conceptual-design support can be valuable for student teams \cite{paper-0046}, but it also shows that usability is not equivalent to verification capability.  Lightweight practice therefore needs an evidence ladder: accessibility and usability are entry criteria, while value realization requires traceable decisions, executed checks, and eventually test closure.

The strongest practice pattern in the included studies combines bounded scope, explicit roles, and recurring review cadence.  The 3ColStar case links Capella/Arcadia modeling with weekly sprints, permissions, subsystem responsibilities, traceability, and preliminary-design outputs \cite{paper-0078}.  Academic project-management frameworks similarly connect model artifacts to PBS/WBS/CBS and review gates, although they do not quantify ROI \cite{paper-0023}.  Such examples suggest that adoption playbooks should include not only modeling diagrams and tool scripts, but also governance: who owns requirements, who approves assumption changes, how power-budget changes propagate, what evidence is required at PDR/CDR-like gates, and which artifacts are frozen or regenerated for test planning.

\subsection{Threats to validity and limitations of this review}
\label{sec:discussion-threats-limitations}

Several threats limit the strength and generality of this discussion.  The search scope was limited primarily to Scopus and IEEE Xplore, with Google Scholar used only for targeted exact-title verification during recovery searches.  Grey literature, internal mission documentation, tool repositories without papers, and industrial reports may contain stronger evidence for MBSE value or model-linked V\&V than the accessible academic literature.  Four potentially eligible reports were not retrieved or remained inaccessible and were excluded from extraction, so the synthesis reflects 40 accessible studies rather than the full set of potentially relevant work.

Extraction and synthesis also have reliability limits.  The review used a predefined protocol, standardized extraction form, run-specific codebook, supporting quotations, quality appraisal, and extraction-confidence ratings, but no independent inter-rater reliability assessment was performed.  Claim-evidence confidence, EPS relevance, V\&V linkage, and Table~5-inspired gap mappings therefore include reviewer judgment.  Controlled-code mapping helped make counts reproducible and reduced the earlier keyword-categorization problem, but it can still simplify nuanced paper content, especially where authors use broad terms such as validation, verification, digital engineering, traceability, or model integration differently.  Emergent codes were retained to keep such observations available for qualitative interpretation.

The evidence base is immature: all 40 studies carried at least one structured limitation flag related to incomplete lifecycle coverage, weak claim--evidence linkage, single-case or limited-scope evidence, no independent validation, and prototype/future-work maturity; 39 lacked physical-test or flight evidence; 38 lacked a baseline or comparator; 35 lacked quantitative MBSE-outcome evaluation; and 34 did not make model assets available.  Repeated publications in the CubeSat System Reference Model lineage also mean that publication count should not be interpreted as independent replication.  Because only two included studies are EPS-specific, EPS conclusions necessarily rely on a mixture of whole-nanosatellite, EPS-adjacent, and power-budget/electrical-architecture studies.

Use of Buruso et al.'s Table~5 gaps is another limitation.  The gaps were used as a positioning and sensitizing frame, not as an independent coding validation instrument \cite{Buruso2026MBSESpace}.  Mapping nanosatellite studies onto those broad gaps is useful for comparison, but the resulting categories are coarse and should be interpreted as evidence-strength descriptors rather than definitive classifications of each study.

